\section{Additional Benchmark Details and Original Results}
\label{app:baselines}
Tables~\ref{tab:queryner-benchmark}--\ref{tab:diagnostic-benchmark} present the main saved benchmark results. Table~\ref{tab:vertical-benchmark} adds the constraint results by vertical: seven fashion, seven grocery, and six general-merchandise queries. CommerceCore is a single served-model evaluation; API means cover three calls per query. The overall mean weights queries equally rather than averaging the three vertical means. Historical notes describe the Sonnet repetitions as identical, but omitting a temperature parameter is not itself proof of deterministic generation.

The six-query API file contains aggregates and aliases, without per-query responses or the provider-side execution script. Its values cannot support paired significance testing or a response-level fairness audit. The GPU-side diagnostic file has parsed predictions and gold labels, permitting rescoring but not inspection of raw generations.

\begin{table}[h]
\centering
\caption{Constraint F1 by vertical: CommerceCore is rescored from archived served-model predictions; API values come from the three-repeat aggregate file. These are development-set results under differing protocols.}
\label{tab:vertical-benchmark}
\begin{tabular}{lrrrr}
\toprule
Stored model / alias & Fashion & Grocery & General & Overall \\
\midrule
CommerceCore, served model & 0.690 & 0.667 & 0.694 & 0.683 \\
\midrule
Claude Haiku 4.5 & 0.803 & 0.448 & 0.267 & 0.518 \\
Claude Sonnet 5 & 0.795 & 0.643 & 0.294 & 0.592 \\
GPT-4o mini / cheap tier & 0.875 & 0.605 & 0.333 & 0.618 \\
GPT-5 / flagship & 0.759 & 0.540 & 0.167 & 0.504 \\
\bottomrule
\end{tabular}
\end{table}

A separate single-run flagship artifact records 0.608 for the Sonnet configuration and 0.542 for GPT-5 on the same development set. These are distinct runs from the three-repeat aggregates. We preserve their files rather than merging unlike measurements. Historical scripts convert API errors into empty predictions or zeros, so end-to-end scores can include request failures.

\paragraph{Preserving the original manuscript's numerical record.}
Table~\ref{tab:original-results} retains the two remaining QueryNER values from the \href{https://github.com/arghya05/commercecore/blob/09c65f3/paper/commercecore_paper.tex}{original manuscript}. A matching result file or prediction trace for these two values has not been located in the repository or the inspected model archive. Their presence here records what the original paper reported; it does not independently verify the measurements or establish comparability with the recorded runs above. The original served-model and open-model constraint results, in contrast, have archive evidence and are included in Table~\ref{tab:constraint-benchmark}.

\begin{table}[htbp]
\centering\small
\caption{Previously reported benchmark values without a matching local measurement artifact. All entries in this table have narrative-only provenance and are excluded from comparative conclusions. Model aliases follow the original paper.}
\label{tab:original-results}
\begin{tabularx}{\linewidth}{@{}p{0.36\linewidth}rX@{}}
\toprule
System / evaluation & Reported F1 & Scope reported in the original manuscript \\
\midrule
GPT flagship / QueryNER & 0.204 & Entity segmentation; sample and execution details unverified. \\
Claude Sonnet 5 / QueryNER & 0.427 & Entity segmentation; sample and execution details unverified. \\
\bottomrule
\end{tabularx}
\end{table}

\section{Training Logs, Checkpoint History, and Reproducibility}
\label{app:training-logs}
\subsection{Complete loss trace and approximate epoch summaries}
The released clean-run log contains one loss record for every optimizer step from 1 to 1,200. At each step, the training script averages the completion-token losses of eight microbatches; the logged value is rounded to four decimal places. It is an average of microbatch losses, not a token-weighted average over an epoch. Figure~\ref{fig:training-loss} plots every logged value and a trailing mean over at most 50 steps. The original CSV is included unchanged in the source package as \path{evidence/training_loss_history.csv} and is available in the repository at \href{https://github.com/arghya05/commercecore/blob/main/train/results_multitask_clean/multitask_training_loss_log.csv}{the complete training-loss log}.

\begin{figure}[htbp]
\centering
\begin{tikzpicture}
\begin{axis}[width=0.95\linewidth,height=5.6cm,xlabel={Optimizer step},ylabel={Logged training loss (log scale)},
 xmin=0,xmax=1200,ymode=log,xtick={0,150,300,450,600,750,900,1050,1200},
 tick label style={font=\small},label style={font=\small},
 legend style={font=\small,at={(0.5,1.03)},anchor=south,draw=none,legend columns=2},grid=major,grid style={black!10}]
\addplot[black!25,thin,no marks] table[x=step,y=loss,col sep=comma]{evidence/training_loss_curve.csv};
\addlegendentry{All 1,200 logged losses}
\addplot[ccBlue,thick,no marks] table[x=step,y=trailing_50_mean,col sep=comma]{evidence/training_loss_curve.csv};
\addlegendentry{Trailing 50-step mean}
\draw[dotted,black] (axis cs:750,0.001)--(axis cs:750,3);
\end{axis}
\end{tikzpicture}
\caption{Training loss from the archived run, including individual-step fluctuations. The dotted line marks selected checkpoint 750. Smoothing summarizes the stored loss values and does not add observations; no validation-loss trace was recorded.}
\label{fig:training-loss}
\end{figure}

One pass through 2,130 examples spans 266.25 optimizer steps at accumulation eight. Table~\ref{tab:epoch-loss} uses boundaries $\lfloor k\cdot2{,}130/8\rfloor$, followed by the final partial pass. Some optimizer steps straddle pass boundaries, so these are approximate epoch windows. Means are arithmetic averages of the rounded logged values; the end-step column also preserves the epoch-boundary loss view used in the project record. Lower training loss alone does not establish improved generalization.

\begin{table}[htbp]
\centering\small
\caption{Loss by approximate training pass, recomputed from all 1,200 records. The final row covers the incomplete fifth pass.}
\label{tab:epoch-loss}
\begin{tabular}{rrrrr}
\toprule
Steps & Passes at end & Logged steps & Mean loss & End-step loss \\
\midrule
\inputtablerows evidence/loss_window_rows.tex
\bottomrule
\end{tabular}
\end{table}

\subsection{Every checkpoint with its exact training-loss record}
\begin{table}[h]
\centering\small
\caption{Complete eight-checkpoint history joined to the exact loss-log step. C-F1 uses 20 constraint development queries; QN-F1 uses the first 30 QueryNER test examples; similarity uses five lexical probes. Elapsed minutes are logged before that step's checkpoint evaluation and saving. Step 750 is selected.}
\label{tab:checkpoint-log}
\begin{tabular}{rrrrrrr}
\toprule
Step & Passes & Loss & C-F1 & QN-F1 & Similarity & Minutes \\
\midrule
\inputtablerows evidence/checkpoint_rows.tex
\bottomrule
\end{tabular}
\end{table}

At step 1,200, the loss log shows 4,385.8 seconds (73.10 minutes), while the run summary records 4,510.6 seconds (75.18 minutes). The latter also includes the final checkpoint evaluation and saving. A step's loss is the eight-microbatch mean for that update, not the mean loss over all preceding training. The table uses the actual checkpoint steps rather than nearby illustrative loss samples.

The local selected adapter configuration agrees with rank 16, alpha 32, dropout 0.05, and the seven projection targets. It identifies PEFT version 0.21.0 and a pod-local base path, but no base revision. This is not a complete environment lock. The training script uses standard AdamW rather than a paged optimizer.

The exact diagnostic prompts are: ``What is the capital of France?'', ``Write a short sentence about the weather today.'', ``What is 12 plus 15?'', ``Name one primary color.'', and ``Complete this sentence: The sun rises in the''. Decoding is greedy with 30 new tokens. Only similarity against the frozen base is logged; assigning a variable for the preceding checkpoint does not create an additional retention measurement.

\begin{figure}[h]
\centering
\begin{tikzpicture}
\begin{axis}[width=0.86\linewidth,height=5.7cm,xlabel={Base-output similarity (five probes)},ylabel={Constraint development F1},
 xmin=0.477,xmax=0.634,ymin=0.665,ymax=0.766,grid=major,grid style={black!10},
 tick label style={font=\small},label style={font=\small},point meta=explicit symbolic,
 nodes near coords, every node near coord/.append style={font=\scriptsize,anchor=south}]
\addplot[only marks,mark=*,ccBlue] table[x=output_similarity,y=constraint_f1,meta=step,col sep=comma]{evidence/checkpoint_history.csv};
\addplot[black!55,dashed,no marks] coordinates {(0.5175,0.7500) (0.6143,0.7250)};
\addplot[only marks,mark=o,mark size=4pt,black,thick] coordinates {(0.6143,0.7250)};
\end{axis}
\end{tikzpicture}
\caption{Constraint-F1/similarity projection with optimizer-step labels. The dashed segment connects the two nondominated points in this two-dimensional projection; it is not an interpolated achievable frontier. The circled point is selected step 750. Including QueryNER F1 produces the larger three-dimensional nondominated set reported in the main text. Similarity is not general-capability accuracy.}
\end{figure}

\section{Result-Level Provenance}
\label{app:provenance}
\paragraph{Public code, models, and data.}
Readers can access the \href{https://github.com/arghya05/commercecore}{GitHub project}, \href{https://huggingface.co/arghya2030/commercecore-qwen3-1.7b}{CommerceCore model card}, \hfartifact{model.safetensors}{merged model weights}, \hfartifact{config.json}{model configuration}, and \hfartifact{runpod_archive/checkpoints_multitask/step_750/adapter_model.safetensors}{selected step-750 adapter}. The upstream resources are the \href{https://huggingface.co/Qwen/Qwen3-1.7B}{Qwen3-1.7B base model}, \href{https://huggingface.co/datasets/bltlab/queryner}{QueryNER dataset card}, and \href{https://aclanthology.org/2024.lrec-main.1178/}{QueryNER paper}. Local dataset versions are linked as the \gitartifact{data/queryner/train.jsonl}{QueryNER training records}, \gitartifact{data/queryner/test.jsonl}{QueryNER test records}, \gitartifact{data/synthetic_training_text.jsonl}{synthetic query targets}, \gitartifact{data/catalog_training_text.jsonl}{catalog targets}, and \gitartifact{data/held_out_eval_set.jsonl}{20-query development set}.

\paragraph{Fixed evidence versions.}
The GitHub evidence links in this appendix resolve to commit \code{fa64ee7}; Hugging Face file links resolve to revision \code{425c4a5}. These fixed versions prevent later edits from silently changing the cited evidence. The \href{https://github.com/arghya05/commercecore/blob/main/paper/evidence/archive/manifest.json}{archive manifest} records the full Hugging Face revision, file URLs, and SHA-256 hashes of the mirrored inputs. Current manuscript and audit links point to the maintained paper folder; the \href{https://github.com/arghya05/commercecore/blob/main/paper/evidence/build_manifest.json}{build manifest} identifies the exact PDF and source-package contents.

\begin{table}[h]
\centering\small
\caption{Clickable evidence map for the reported numbers. Offline rescoring verifies the stored predictions and arithmetic without rerunning model inference.}
\begin{tabularx}{\linewidth}{@{}>{\raggedright\arraybackslash}p{0.25\linewidth}>{\raggedright\arraybackslash}X@{}}
\toprule
Result & Artifact and verification boundary \\
\midrule
Mixture and seed-42 sample & \gitartifact{train/build_multitask_mixture.py}{Mixture builder}; \hfartifact{runpod_archive/data/multitask_train_mixture.jsonl}{archived 2,130-row mixture}. Reconstructed records match in order. \\
75.18 minutes; \$0.9021 & \gitartifact{train/results_multitask_clean/multitask_run_summary.json}{Run summary}. Cost arithmetic verified; not a provider invoice. \\
Checkpoint trajectories & \gitartifact{train/results_multitask_clean/multitask_checkpoint_eval_log.csv}{Task-F1 CSV}; \gitartifact{train/results_multitask_clean/multitask_retention_log.csv}{similarity CSV}; \gitartifact{train/results_multitask_clean/MODEL_SELECTION_DECISION.md}{selection note}. Eight checkpoints. \\
Loss curve and epoch means & \gitartifact{train/results_multitask_clean/multitask_training_loss_log.csv}{All 1,200 loss records}; \hfartifact{runpod_archive/logs/train_multitask_v2.log}{training console log}. Summaries and exact-step joins recomputed. \\
Served F1 and category scores & \hfartifact{runpod_archive/results/complete_inference_test_results.json}{All 20 served predictions and labels}; \hfartifact{runpod_archive/scripts/complete_inference_test.py}{evaluation script}. F1 0.683 and vertical means reproduced. \\
Open-model development F1 & \hfartifact{runpod_archive/results/open_model_comparison_results.json}{Stored aggregates}; \hfartifact{runpod_archive/scripts/open_model_comparison.py}{evaluation script}. 0.050 and 0.100; predictions absent. \\
Full QueryNER F1 and CI & \gitartifact{eval/full_queryner_eval_result.json}{Result JSON}; \gitartifact{eval/full_queryner_eval.py}{scoring script}. Historical per-example scores absent. \\
API QueryNER results & \gitartifact{eval/queryner_baseline_results.json}{Result JSON}; \gitartifact{eval/queryner_baseline_comparison.py}{prompt and evaluation script}. Fifty queries, restricted labels. \\
Repeated API constraints & \gitartifact{eval/rigorous_baseline_results.json}{Result JSON}; \gitartifact{eval/rigorous_baseline_comparison.py}{evaluation script}. Pooled-repeat CIs excluded from inference. \\
Single-run flagship results & \gitartifact{eval/flagship_tier_comparison_results.json}{Result JSON}; \gitartifact{eval/flagship_tier_comparison.py}{evaluation script}. Kept separate from three-repeat runs. \\
Six-query local/API results & \gitartifact{eval/fresh_headtohead_gpu_results.json}{Local-model predictions}; \gitartifact{eval/fresh_headtohead_frontier_results.json}{API aggregates}; \gitartifact{eval/fresh_head_to_head_comparison.py}{query definitions}. Saved local predictions rescored. \\
Generation cost & \gitartifact{data/generation_cost_ledger.json}{Generation ledger}. Sum of 1,630 call estimates verified. \\
Offline analyses and hashes & \href{https://github.com/arghya05/commercecore/blob/main/paper/evidence/audit_results.json}{Audit results}; \href{https://github.com/arghya05/commercecore/blob/main/paper/audit_evidence.py}{audit code}. Includes archived-mixture checks and served-model rescoring. \\
\bottomrule
\end{tabularx}
\end{table}

Run the offline audit from the repository root with \code{python3 paper/audit\_evidence.py}. Normalization uses Unicode case folding and whitespace collapse, without punctuation removal or semantic deduplication. The \href{https://github.com/arghya05/commercecore/blob/main/paper/evidence/reconstructed_queryner_train_ids.txt}{reconstructed QueryNER sample IDs} are saved explicitly; the reconstructed prompt/completion sequence matches the archived mixed training file. The hash manifest identifies exactly which data and scripts were audited.

The implementation links are the \gitartifact{train/train_multitask.py}{training loop}, \gitartifact{data/generate_scaled_query_training_data.py}{query generator}, \gitartifact{validators/mechanical_validator.py}{validator}, and \hfartifact{runpod_archive/scripts/serve_inference.py}{archived serving parser}. The eight-query \hfartifact{runpod_archive/results/full_inference_test_results.json}{JSON-format smoke test} stores eight valid JSON outputs, which are independently parseable; it lacks gold constraints and is not an extraction-accuracy benchmark. The maintained \href{https://github.com/arghya05/commercecore/blob/main/paper/CommerceCore_Paper.pdf}{paper PDF}, \href{https://github.com/arghya05/commercecore/blob/main/paper/commercecore_paper.tex}{LaTeX source}, and \href{https://github.com/arghya05/commercecore/blob/main/paper/CommerceCore_arXiv_Source.zip}{self-contained source ZIP} are linked for inspection and rebuilding.

Table~\ref{tab:original-results} preserves the original manuscript's additional benchmark values with their evidence status. A narrative claim or executable script alone is not a saved measurement. A comment also records a 313/450 generic-template issue in an earlier batch, but that batch is unavailable for recounting; no measured before/after quality improvement is inferred.

The original planning documents describe broader retrieval, tool execution, and merchant-transfer objectives. A plan is not evidence of a completed experiment. This paper reports the implemented extraction study only.

\section{Protocol for a Confirmatory Experimental Extension}
\label{app:future}
This section specifies additional experiments motivated by the study, not completed results. The intended questions are the effect of task specialization, the contribution of each data source, robustness to linguistic and schema shifts, and accuracy--latency--cost trade-offs.

\paragraph{Data integrity and independent evaluation.}
Freeze an annotation contract for supported fields, negation, unit conversion, and unsupported requests. Serialize the validated target itself and audit completeness as well as span binding. Group repeated scenarios and paraphrases before splitting. Use development data for checkpoint and prompt selection. In a new training cycle, use the QueryNER development split rather than test examples. Do not relabel historical selection data as unseen. Build a separate constraint test set with documented provenance and label validation.

\paragraph{Robustness coverage.}
A planned 300--1,000-query benchmark should stratify one, two, three, and four-or-more constraints; negation; units; implicit values; paraphrases; and typos. Counts alone do not establish independence or label validity. Separate linguistically perturbed scenarios from merchant/schema transfer and include genuinely distinct scenarios or sources where permitted. Determine sample size from a prespecified precision or power target, not an observed favorable difference. Labels need independent validation; mechanically generated cases should be identified as synthetic stress tests.

\paragraph{Specialist baselines and ablations.}
Evaluate the unfine-tuned Qwen3-1.7B backbone, QueryNER-style BERT and domain-adapted BERT, a modern encoder tagger, a small sequence-to-sequence extractor, deterministic rules, API models, and CommerceCore on identical records. Supply the full ontology, appropriate templates, and sufficient generation budgets. Entity-only encoders should be compared on the entity task unless a documented normalization/extraction head makes constraint outputs comparable. Continued pretraining entails additional data and compute that must be reported.

Match presentation counts, token budgets, tuning allowance, and decoding for data-source and task-sharing ablations. ``Real only'' contains entity supervision but no synthetic constraint targets; that setting changes task coverage as well as data origin and should be interpreted accordingly. Table~\ref{tab:ablationplan} defines the scope without inventing scores.

\begin{table}[h]
\centering\small
\caption{Planned ablations; no results have been produced for this table. C = synthetic query constraints; E = human entity labels; A = synthetic catalog attributes. Historical checkpoint comparisons are not substitutes for these controls.}
\label{tab:ablationplan}
\begin{tabularx}{\linewidth}{@{}lclX@{}}
\toprule
Configuration & Sources & Status & Interpretation \\
\midrule
Real only & E & Not run & Entity training; constraint task lacks direct targets. \\
Synthetic only & C+A & Not run & Removes human entity supervision. \\
Real + query synthetic & E+C & Not run & Removes catalog task from the mixture. \\
No phrase gate & E+C+A & Not run & Requires archived pre-gate data and matched sampling. \\
No QueryNER multitask & C+A & Not run & Same data-source condition as synthetic only. \\
Corrected full setup & E+C+A & Not run & Repairs label integrity; historical results cannot be reused. \\
\bottomrule
\end{tabularx}
\end{table}

\paragraph{Inference and uncertainty.}
Run multiple training seeds and archive raw generations, parsed outputs, failure flags, per-query scores, model revisions, and settings. Prespecify a primary endpoint and paired query-level uncertainty; report exact-set correctness, precision/recall, and results by phenomenon and vertical. Use a scored general-capability suite for retention claims. Evaluate batch sizes 1, 8, 32, and 64 where memory permits, including sustained throughput, p50/p95 latency, full device memory, startup, utilization, and comparable API request costs. Record failures rather than omitting infeasible operating points.
